% from ega0-17-fr.tex
\oldpage[0\textsubscript{IV-1}]{42}

\begin{corollary}[17.1.8]
\label{0.17.1.8-fr}
Soient $A$ un anneau local noethérien, $\mathfrak{m}$ son idéal maximal, $t$ un
élément de $\mathfrak{m}$. Les conditions suivantes sont équivalentes~:
\begin{enumerate}
  \item[\rm a)] $A/tA$ est régulier et $t$ est non diviseur de zéro dans $A$.
  \item[\rm b)] $A$ est régulier et $t\notin\mathfrak{m}^2$.
\end{enumerate}
\end{corollary}

\begin{proof}
C'est le cas particulier $r=1$ de (17.1.7) (compte tenu de ce qu'un élément de
$\mathfrak{m}$ non diviseur de zéro fait partie d'un système de paramètres
(16.4.1)).
\end{proof}

\begin{corollary}[17.1.9]
\label{0.17.1.9-fr}
Soient $A$ un anneau local régulier, $\mathfrak{m}$ son idéal maximal,
$\mathfrak{J}$ un idéal de $A$ contenu dans $\mathfrak{m}$. Les conditions
suivantes sont équivalentes~:
\begin{enumerate}
  \item[\rm a)] L'anneau $A/\mathfrak{J}$ est régulier.
  \item[\rm b)] $\mathfrak{J}$ est engendré par une suite
  $(x_i)_{1\leq i\leq r}$ faisant partie d'un système régulier de paramètres
  de $A$.
\end{enumerate}
\end{corollary}

\begin{proof}
On a déjà vu (17.1.7) que b) entraîne a). Inversement, supposons que
$A/\mathfrak{J}$ soit régulier (ce qui entraîne que $\mathfrak{J}$ est premier)
et soient $n=\dim(A)$, $n-r=\dim(A/\mathfrak{J})$~; avec les notations de
(17.1.7), la suite exacte (17.1.7.1) donne
$\operatorname{rg}_k((\mathfrak{m}^2+\mathfrak{J})/\mathfrak{m}^2)=r$, donc il
existe une suite $(x_i)_{1\leq i\leq r}$ d'éléments de $\mathfrak{J}$ faisant
partie d'un système régulier de paramètres de $A$. Posons
$\mathfrak{J}'=x_1A+\cdots+x_rA$~; il résulte de (17.1.7) que $\mathfrak{J}'$
est un idéal premier de $A$ et que $A/\mathfrak{J}'$ est régulier et de
dimension $n-r$~; mais comme
$A/\mathfrak{J}=(A/\mathfrak{J}')/(\mathfrak{J}/\mathfrak{J}')$ et que
$\mathfrak{J}/\mathfrak{J}'$ est premier dans l'anneau intègre
$A/\mathfrak{J}'$, les dimensions de $A/\mathfrak{J}$ et de
$A/\mathfrak{J}'$ ne peuvent être égales que si
$\mathfrak{J}=\mathfrak{J}'$ (16.1.2.2).
\end{proof}

\subsection{Rappels sur la dimension projective et la dimension injective des
modules}
\label{subsection:0.17.2-fr}

\begin{env}[17.2.1]
\label{0.17.2.1-fr}
Soient $A$ un anneau, $M$ un $A$-module. Rappelons (M, VI, 2) qu'on appelle
\emph{dimension projective} (resp. \emph{dimension injective}) de $M$ et que
l'on note $\operatorname{dim.proj}(M)$ ou
$\operatorname{dim.proj}_A(M)$ (resp. $\operatorname{dim.inj}(M)$ ou
$\operatorname{dim.inj}_A(M)$) le plus petit $n$ (entier ou égal à $+\infty$)
tel qu'il existe une résolution projective gauche (resp. une résolution
injective droite) de $M$ qui soit \emph{de longueur $n$}. Il revient au même
(\emph{loc. cit.}) de dire que $n$ est le plus petit nombre tel que pour tout
$A$-module $N$, on ait $\operatorname{Ext}^i_A(M,N)=0$ (resp.
$\operatorname{Ext}^i_A(N,M)=0$) pour tout $i>n$, ou seulement pour $i=n+1$.
Pour tout \emph{facteur direct} $M'$ de $M$, on a donc
$\operatorname{dim.proj}(M')\leq\operatorname{dim.proj}(M)$ puisque
$\operatorname{Ext}^i_A(M',N)$ est facteur direct de
$\operatorname{Ext}^i_A(M,N)$~; de même
$\operatorname{dim.inj}(M')\leq\operatorname{dim.inj}(M)$. Une condition
équivalente à $\operatorname{dim.proj}(M)=n$ (resp.
$\operatorname{dim.inj}(M)=n$) est que $n$ est le plus petit nombre tel que,
pour toute suite exacte
\[
  0\longrightarrow R\longrightarrow P_{n-1}\longrightarrow\cdots
  \longrightarrow P_0\longrightarrow M\longrightarrow0
\]
(resp.
$0\longrightarrow M\longrightarrow Q_0\longrightarrow\cdots
\longrightarrow Q_{n-1}\longrightarrow C\longrightarrow0$), où tous les
$P_i$ pour $0\leq i<n$ sont projectifs (resp. tous les $Q_i$ pour $0\leq i<n$
sont injectifs), alors $R$ est projectif (resp. $C$ est injectif).
\end{env}

\begin{remarks}[17.2.2]
\label{0.17.2.2-fr}
\begin{enumerate}
  \item[\rm (i)] Dire que $M$ est un $A$-module \emph{projectif} (resp.
  \emph{injectif}) équivaut donc à dire que
  $\operatorname{dim.proj}(M)=0$ (resp. $\operatorname{dim.inj}(M)=0$).
  \item[\rm (ii)] Soit $T$ un \emph{foncteur covariant} additif de la catégorie
  des $A$-modules dans une catégorie abélienne. Il résulte aussitôt des
  définitions de (17.2.1) et de celle des
  